\documentclass[11pt]{article}
\usepackage[margin=0.9in]{geometry}
\usepackage{amsmath,amssymb}
\usepackage{enumitem}
\usepackage{tcolorbox}
\usepackage{xcolor}
\usepackage{booktabs}
\usepackage{array}
\usepackage{fancyhdr}
\tcbuselibrary{skins,breakable}

\definecolor{boxbg}{gray}{0.94}
\definecolor{boxline}{gray}{0.55}
\definecolor{keybg}{gray}{0.97}

\pagestyle{fancy}
\fancyhf{}
\renewcommand{\headrulewidth}{0pt}
\fancyfoot[C]{\small\color{gray} Calculus 1 \quad$\cdot$\quad The Frost Night Water Decision \quad$\cdot$\quad Case Study}
\fancyfoot[R]{\small\color{gray}\thepage}

\newtcolorbox{databox}{enhanced, breakable, colback=boxbg, colframe=boxline,
  boxrule=0.6pt, arc=3pt, left=10pt, right=10pt, top=8pt, bottom=8pt}
\newtcolorbox{keybox}{enhanced, breakable, colback=keybg, colframe=boxline,
  boxrule=0.8pt, arc=3pt, left=12pt, right=12pt, top=10pt, bottom=10pt}

\newcommand{\question}[1]{\medskip\noindent\textbf{#1}\par\smallskip}
\newcommand{\hint}[1]{{\small\color{gray}\textit{Hint: #1}}\par}
\newcommand{\skill}[1]{{\small\textit{Skill: #1}}\par}

\begin{document}
{\small NAME: \underline{\hspace{2.6in}}}\par\vspace{6pt}
\begin{center}
{\LARGE\bfseries The Frost Night Water Decision}\\[3pt]
\rule{2.2in}{0.8pt}\\[5pt]
{\itshape A frost-night supply decision for small groups \quad$\bullet$\quad Calculus 1}
\end{center}

A hard freeze is forecast for sunrise tomorrow at Cold Hollow Orchards, and operations lead Marisol Vega must run frost sprinklers over the 12-acre block of early blossoms through the coldest morning hours. The sprinklers draw from a holding tank that was drained for cleaning and refills overnight from the orchard's well, starting at midnight. Working from last week's metered test run, the district hydrologist models the overnight inflow at $W(t) = 200t - 25t^2$ gallons per hour, where $t$ is hours after midnight, valid through 6 a.m. The sprinklers need at least 2,000 gallons in the tank by 6 a.m. A local hauler can deliver 800 gallons for \$190, but bookings close at 11 p.m. It is 10 p.m. Marisol must decide: trust the well alone, or book the truck.

\begin{databox}
\begin{center}
\begin{tabular}{@{}l p{3.9in}@{}}
\toprule
Refill window & midnight ($t = 0$) to 6 a.m. ($t = 6$); the tank starts empty \\
Modeled well inflow & $W(t) = 200t - 25t^2$ gallons per hour, valid $0 \le t \le 6$ \\
Frost requirement & at least 2,000 gallons in the tank by 6 a.m. \\
Filter skid limiter & halts the refill if the flow ever exceeds 450 gallons per hour \\
Backup hauler & one 800-gallon delivery, \$190, booking closes at 11 p.m. \\
Holding tank & capacity 3,000 gallons; serves the 12-acre blossom block \\
\bottomrule
\end{tabular}
\end{center}
\smallskip
\noindent\textbf{Pump-house reference card} (district Metering Handbook): when a flow varies during a metering window, it spends time below its highest reading, so its average flow for the window (total gallons divided by hours of flow) is less than that highest reading. Estimating total gallons as ``highest reading $\times$ hours'' therefore overstates the true total whenever the flow varies.
\end{databox}

\question{1. Check the Overnight Peak Flow}
At what time overnight does the well deliver water fastest, and at what rate in gallons per hour? State whether the filter skid's limiter will interrupt the refill at any point in the window.
\hint{A night's fastest flow can hide at the edges of the window: compare the overnight high point against the flows at midnight and at 6 a.m.}

\question{2. Total the Night's Refill}
How many gallons does the well actually add to the tank between midnight and 6 a.m.? Compare your total to the frost requirement: does the well alone get Marisol to 2,000 gallons by sunrise, and if not, by how much does it fall short?
\hint{The flow is different every hour, so no single hour's reading can speak for the whole night.}

\question{3. Settle the Hauler Dispute}
Dev, the assistant manager, wants to save the \$190: ``The well hits 400 gallons an hour, and we run it for six hours. That is 2,400 gallons, comfortably past 2,000. Cancel the truck.'' Compute Dev's estimate for yourselves, then discuss with your group: whose total should Marisol trust, Dev's or yours from the previous part? Use the pump-house reference card to explain which side it supports. Then decide whether to book the delivery.
\hint{How much of the night does the well really spend at its very best rate?}

\bigskip
\noindent\textbf{Your Recommendation.} As a group, write 2--3 sentences to Marisol: tell her whether to book the hauler before 11 p.m., citing the totals that settle the dispute and what the \$190 actually buys her.

\vspace{10pt}
\noindent\rule{\textwidth}{0.4pt}\par\vspace{16pt}
\noindent\rule{\textwidth}{0.4pt}\par\vspace{16pt}
\noindent\rule{\textwidth}{0.4pt}

\newpage
\section*{Instructor Material}

\subsection*{Analyst's Checklist (hand out only if a group is stuck)}
\begin{enumerate}[leftmargin=*, itemsep=2pt]
\item Write down the well's flow at midnight and at 6 a.m. using the model.
\item Find the time when the flow stops rising and starts falling, and compute the flow at that moment.
\item Compare all three flows; the largest is the night's peak. Check it against the 450 gal/hr limiter.
\item Find the total gallons the changing flow delivers across the whole six-hour window, not a one-hour figure scaled up by six.
\item Compare that total to the 2,000-gallon requirement, then test what one 800-gallon delivery changes.
\item For Dev's claim, ask: what would the flow have to look like all night long for ``peak times hours'' to be exact?
\end{enumerate}

\subsection*{Answer Key}
\begin{keybox}
\textbf{Part 1: peak flow.}\par
\skill{Locating an absolute maximum on a closed interval: interior critical point plus endpoint comparison (OpenStax Calculus Vol.~1, Sec.~4.3)}
$W'(t) = 200 - 50t$; setting $W'(t) = 0$ gives $t = 4$, an interior point of $[0,6]$.\\
Compare candidates: $W(0) = 0$, \quad $W(4) = 800 - 400 = 400$, \quad $W(6) = 1200 - 900 = 300$.\\
$\boxed{\text{Peak flow } 400 \text{ gal/hr at } t = 4 \text{ (4 a.m.)}}$\quad Since $400 < 450$, the limiter never trips and the refill runs uninterrupted.

\medskip
\textbf{Part 2: the true total.}\par
\skill{Net change from a rate function: the total added is the definite integral of the inflow rate (Sec.~5.4), evaluated by antiderivative subtraction (Sec.~5.3)}
$\displaystyle \int_0^6 \left(200t - 25t^2\right)dt = \left[100t^2 - \tfrac{25}{3}t^3\right]_0^6 = (3600 - 1800) - 0 = 1800.$\\
$\boxed{\text{The well adds } 1800 \text{ gallons}}$\quad Since $1800 < 2000$, the well alone leaves Marisol \textbf{200 gallons short} of the frost requirement.

\medskip
\textbf{Part 3: the dispute and the decision.}\par
\skill{Average value versus peak value of a varying rate; error analysis of a constant-rate assumption (Sec.~5.3 and the printed reference card)}
Dev's estimate: $400 \times 6 = 2400$ gallons, which would clear 2,000 and cancel the truck.\\
But the well flows at 400 gal/hr only at the single instant $t = 4$; at every other moment it is slower. The night's average flow is $1800/6 = 300$ gal/hr, below the 400 peak, exactly as the reference card says.\\
Both totals are clean and the 2,000-gallon requirement sits strictly between them ($1800 < 2000 < 2400$): Dev's shortcut does not just miss, it flips the decision.\\
$\boxed{\text{Book the delivery: } 1800 + 800 = 2600 \ge 2000}$

\medskip
\textbf{Verdict.} Book the hauler before 11 p.m. The well alone delivers 1,800 gallons, 200 short of the requirement, while Dev's peak-times-hours figure of 2,400 overstates the refill by 600 gallons because the well spends all night except one instant below its peak. The \$190 delivery lifts the tank to 2,600 gallons, a 600-gallon margin, still comfortably under the 3,000-gallon capacity.
\end{keybox}

\subsection*{Alignment}
\begin{itemize}[leftmargin=*, itemsep=2pt]
\item \textbf{Source problem objective:} determine a rate function's maximum value via critical-point analysis, compute the net accumulated total via the net change theorem, compare it to a decision threshold, and explain why estimating that total from the peak rate overstates it.
\item \textbf{PRIMARY section:} OpenStax \textit{Calculus Volume 1} (Strang/Herman), Sec.~5.4, ``Integration Formulas and the Net Change Theorem'', the rate-to-total conversion behind Parts 2 and 3.
\item \textbf{SUPPORTING sections:} Sec.~4.3, ``Maxima and Minima'', the critical-number and closed-interval comparison behind Part 1; Sec.~5.3, ``The Fundamental Theorem of Calculus'', the antiderivative evaluation used in Part 2 and the average-value idea behind Part 3.
\item \textbf{Section objectives exercised:} ``Use the net change theorem to solve applied problems'' (5.4); ``Describe how to use critical points to locate absolute extrema over a closed interval'' (4.3); ``Use the Fundamental Theorem of Calculus, Part 2, to evaluate definite integrals'' (5.3).
\item \textbf{Reference data:} the pump-house reference card prints the one fact the cited sections do not state in the exact form needed (a varying flow's average is strictly below its peak, so peak times duration overstates the total), per the mapping's missing-concepts entry; students need no other outside facts.
\item \textbf{What makes this case distinct, and distractors:} the 2,000-gallon requirement sits strictly between the honest total (1,800) and the peak-times-hours estimate (2,400), so the planted misconception flips the booking decision rather than merely perturbing a number. Inert distractors: the 3,000-gallon tank capacity and the 12-acre block size.
\item \textbf{Attribution:} adapted from \textit{Calculus Volume 1} by Gilbert Strang and Edwin ``Jed'' Herman, OpenStax, used under CC BY-NC-SA 4.0. Access for free at openstax.org.
\end{itemize}

% VERIFICATION
% Part 1: W'(t) = 200 - 50t; 200 - 50t = 0 -> t = 4, interior to [0,6]. W(0) = 0; W(4) = 200*4 - 25*16 = 800 - 400 = 400; W(6) = 200*6 - 25*36 = 1200 - 900 = 300. Max of {0, 400, 300} = 400 gal/hr at t = 4 (4 a.m.). Limiter check: 400 < 450, never trips.
% Part 2: antiderivative F(t) = 100t^2 - (25/3)t^3. F(6) = 100*36 - (25/3)*216 = 3600 - 25*72 = 3600 - 1800 = 1800; F(0) = 0. Total = 1800 gal. Requirement: 1800 < 2000, shortfall = 2000 - 1800 = 200 gal.
% Part 3: Dev's estimate 400*6 = 2400 gal; 2400 >= 2000 would cancel the truck. True total 1800 < 2000 requires it; threshold 2000 lies between 1800 and 2400, so the error flips the decision; gap = 2400 - 1800 = 600 gal. Average flow = 1800/6 = 300 gal/hr < 400 peak, consistent with the reference card. With the delivery: 1800 + 800 = 2600 >= 2000 (margin 600), and 2600 <= 3000 capacity, no overflow.

% ATTRIBUTION
% Calculus Volume 1, by Gilbert Strang and Edwin "Jed" Herman, OpenStax. License: CC BY-NC-SA 4.0. Required attribution line per the sections' metadata: "Access for free at openstax.org."

% DERIVED PLAYBOOK (review and promote to library if good)
% P1 SCOPE FENCE: in play: locating an absolute maximum of a rate function on a closed interval via an interior critical number plus endpoint comparison (Sec 4.3); converting a time-varying rate into an accumulated total via the net change theorem with antiderivative evaluation (Secs 5.4 and 5.3); comparing a computed total to a stated threshold; peak-versus-average reasoning supplied as printed reference data. Banned, appearing in neither the cited sections nor the printed data: constrained optimization, related rates, Riemann-sum estimation, substitution, the second-derivative test (the 4.3 extract teaches the closed-interval comparison only), and any average-value machinery beyond the printed card's fact.
% P2 GRAIN SIZE: the source problem's four parts map onto compute (peak), compute + interpret (total vs threshold), and judge (peak-times-duration error). Realized as three interlocking questions: Q1 peak (feeds the trap arithmetic in Q3), Q2 total vs requirement (the decision quantity), Q3 the trap as a teammate's claim plus the booking decision; the group recommendation needs all three. Two fine-grained concepts interlocking fits the 15-20 minute budget.
% P3 FULL STRENGTH: the degenerate version is a monotone or constant rate, where the peak sits at an endpoint and rate-times-time is exact. Full strength here: an interior maximum that must beat BOTH endpoints (0 and 300 are genuine candidates), and a genuinely varying rate whose integral (1800) differs materially from peak times duration (2400).
% P4 CLEAN NUMBERS: coefficients chosen so the critical point is an integer hour, every rate evaluation is a multiple of 10, and the antiderivative evaluation lands on integers; denominator-3 terms are allowed only when the evaluation cancels them, as in (25/3)(216) = 1800. Constants fixed: window [0,6] hours, requirement 2000 gal, hauler 800 gal at $190, limiter 450 gal/hr. Fresh instance: no data value or answer reused from problem.txt or verified_answer.txt (source instance: 300t - 30t^2 on [0,8], peak 750 at t = 5, total 4480, threshold 5000, estimate 6000, average 560).
% P5 THE TRAP: candidate misconceptions for these skills: (a) peak rate times duration equals the total, treating a varying rate as constant at its maximum; (b) reading the rate value W(6) as the accumulated total; (c) assuming the fastest flow occurs at the end of the window. Chosen: (a), the source problem's own part (iv) misconception, because it has decision stakes: wrong path 400 x 6 = 2400 (clean), right path 1800 (clean), requirement 2000 strictly between, so the trap flips the booking decision.
% P6 AUTHENTIC USERS: orchard and vineyard frost-protection managers, irrigation-district and municipal water operators, reservoir and stormwater engineers, fuel-depot and process-plant schedulers. The chosen domain, a frost-night refill against a requirement and a booking deadline, draws on the source's water-storage context but is a fresh decision.
% P7 INTERPRETATION DEMAND: Q1's peak must be read as a clock time (4 a.m.) and checked against the 450 gal/hr limiter (the refill runs uninterrupted); Q2's total must be turned into a shortfall statement against the 2000-gal requirement (200 short); Q3's two totals must be translated into the booking decision and the $190 stake, with the average flow (300) versus the peak (400) explaining who is right.
% P8 REFERENCE DATA: lo_mapping's missing_concepts entry (the average value of a continuous, non-constant rate over a window is strictly below its peak, so peak times duration overstates the true accumulated total) is printed in the student data block as the pump-house reference card, framed as a Metering Handbook fact with no technique named. Nothing further is needed beyond what the cited sections teach.
\end{document}
